\section*{Ethical Framework and IRB Review Process} 

This dissertation employs collaborative autoethnography and practice-based research methodologies that position participants as collaborators and co-researchers (in a sense) rather than as your typical "subjects of study." The research centers the PI's (mine) own speculocultural technopoietic practice, with collaborative elements functioning as shared knowledge creation and validation rather than traditional data collection from human subjects study. 

\subsection*{IRB Review Process and Timeline} 

In an abundance of caution and commitment to ethical research practices, an IRB protocol seeking exemption determination was created on October 20, 2025 and submitted on October 26, 2025 to the Drexel University Institutional Review Board (IRB). The protocol detailed the collaborative autoethnographic methodology and the nature of co-creative participation. 

On October 31, 2025, the IRB indicated \emph{rejection} of the initial protocol without requesting revisions or providing explanation for the determination (documentation included in Appendix \ref{apdx:irb-notice}). Following this notification, I sent an inquiry to the IRB office on November 10, 2025 (after stumbling upon the initial rejection in the "Inbox" section of COEUS Lite) requesting clarification on the grounds for non-approval and guidance for potential revisions.

After nearly a month without response, I sent a second inquiry on December 1, 2025. The IRB subsequently responded with specific feedback regarding required modifications to the protocol. My advisor and I made the requested corrections and resubmitted the revised protocol on December 2, 2025. 

\textbf{Current Status}: As of this proposal submission, the revised IRB protocol is under active review by the Drexel University IRB. No data collection or research activities involving human participants will commence until full IRB approval or exemption determination is obtained.

\subsection*{Unique Considerations for Practice-Based Research}

This research presents methodological characteristics that complicate traditional IRB frameworks, which are primarily designed for experimental or observational social science research. These being it is primarily and Autoethnographic focus, collaborative, and presents shared authorship and ownership with other human collaborators.These characteristics have been communicated to the IRB through the protocol submission materials. The IRB review process will determine the appropriate level of oversight required for this research design.

\subsection*{Ethical Commitments Regardless of IRB Determination}

Irrespective of the final IRB determination, this research maintains rigorous ethical standards grounded in established principles of collaborative and practice-based research:

\begin{description}[noitemsep] 
	\item[Informed Collaboration] All collaborators are fully informed about the nature of the research, how their contributions may be documented or represented, and their rights regarding participation and attribution.
	
	\item[Voluntary and Revocable Participation] Participation is entirely voluntary. Collaborators may withdraw at any time, and any materials or contributions can be removed from the research upon request.
	
	\item[Attribution and Consent] Collaborators will be consulted regarding how they wish to be credited (full attribution, pseudonym, or anonymity). Written consent for any use of names, images, or creative contributions will be obtained.
	
	\item[Cultural Protocols and Respect] Research practices honor the cultural epistemologies, community protocols, and intellectual traditions of all participants. This includes recognition that some knowledge may not be appropriate for public dissemination.
	
	\item[Reciprocity and Mutual Benefit] Collaborative relationships are structured to ensure mutual benefit. Technical knowledge, tools, and creative outputs are shared. Collaboration is not extractive.
	
	\item[Transparency and Accountability] The PI maintains transparency about research purposes and processes. Collaborators have access to materials and can contest representations.
	
	\item[Data Sovereignty] Collaborators retain ownership of their creative contributions. Any recordings, code, or materials created collaboratively will be shared with creators and used only with permission.
\end{description}

\subsection*{Due Diligence in Ethical Research Preparation} 

The PI completed comprehensive training in research ethics: 

\begin{itemize}[noitemsep]
	\item \textbf{CITI Program}: Human Subjects Research training completed in Social \& Behavioral Research branch (Certificates in Appendix \ref{apdx:citi}) 
	\item \textbf{Responsible Conduct of Research}: RCRG 600 completed (Winter Quarter 2023-2024) 
	\item \textbf{IRB Protocol Development}: Prepared comprehensive IRB submission materials (Appendices \ref{apdx:protocol}, \ref{apdx:consent}, \ref{apdx:intent}) demonstrating thorough consideration of participant protections
	\item \textbf{Responsive Revision}: Incorporated IRB feedback and resubmitted revised protocol addressing identified concerns
\end{itemize} 

All IRB submission materials, correspondence, and revision documentation are included in the appendices to demonstrate the thoroughness of the ethical preparation and provide transparency regarding the IRB review process. 

\subsection*{Research Timeline Contingencies}

The anticipated project timeline (Chapter \ref{ch:timeline}) accounts for the IRB approval process. The research design includes the following contingency approaches:

\begin{enumerate}[noitemsep]
	\item \textbf{If IRB Approval/Exemption Obtained}: Proceed with collaborative autoethnographic practice as designed, with full documentation of collaborative sessions and co-creative processes.
	
	\item \textbf{If IRB Review Extended}: The early phases of technical development, infrastructure building, and solo autoethnographic practice can proceed without collaborative participants while awaiting IRB determination.
	
	\item \textbf{If Additional Revisions Required}: Work with IRB to address concerns while maintaining the integrity of the research design and epistemological framework.
\end{enumerate}
